\documentclass[10pt,a4paper]{amsart}
\usepackage[margin=19mm]{geometry}
\usepackage{amsmath,amssymb,mathtools}
\usepackage[scr=rsfs,cal=euler]{mathalfa}
\usepackage{xfrac}
\usepackage[unicode,hidelinks,bookmarks=false]{hyperref}
\newcommand{\ie}{i.e.\ }
\newcommand{\cf}{cf.\ }
\newcommand{\resp}{respectively\ }
\newcommand{\Subsection}{\subsection}
\providecommand{\nobreakdash}{\nobreak\hbox{-}}
\setlength{\emergencystretch}{2em}
\multlinegap0pt
\usepackage{bm}
\usepackage{bbm}
\usepackage{dsfont}
\usepackage{xspace}
\usepackage{cancel}
\usepackage{mathtools}
\usepackage{amsthm}
\makeatletter
\@ifundefined{theorem}{\newtheorem{theorem}{Theorem}}{}
\@ifundefined{lemma}{\newtheorem{lemma}[theorem]{Lemma}}{}
\makeatother
\title[Weyl group symmetries of the toric variety associated with W]{Weyl group symmetries of the toric variety associated with Weyl chambers}
\author{Tao Gui, Hongsheng Hu, Minhua Liu}
\date{Source: arXiv:2410.13617v2 (CC BY 4.0)}
\begin{document}
\maketitle

\noindent\emph{Source and licence.} Weyl group symmetries of the toric variety associated with Weyl chambers. Version arXiv:2410.13617v2 (CC BY 4.0), distributed under the Creative Commons Attribution 4.0 International licence, as recorded in the arXiv licence metadata for this exact version and shown on the abstract page https://arxiv.org/abs/2410.13617v2. Full rights notice: licenses/arxiv-2410.13617-CC-BY-4.0.txt.\par\smallskip

\noindent\textit{Editorial scope.} Counted unit: the Lemma whose source label is \texttt{lem-T-coh-gen} in the article (source0.tex lines 470--473, source label \texttt{lem-T-coh-gen}), delivered together with the article's own complete original proof of it (source0.tex lines 475--522, one \texttt{proof} environment of 48 lines). No mathematics is written, repaired or re-derived: statement and proof are the article's own text line for line. Required context retained uncounted: lem-stab (lines 285--291); lem-stab-2 (lines 287--290); lem-facet (lines 337--343); lem-facet-1 (lines 339--342). Every definition, display and label of those lines is retained unchanged. Omitted from this package and not redistributed: 1--284, 291--336, 343--469, 474--474, 523--622. Nothing inside a retained excerpt is omitted or altered: every retained body is delivered exactly as the article's lines are written, so no omission receipt of any kind accompanies this package. Citations: the retained excerpts carry no citation command at all, so no bibliography entry of the article is transcribed and no reference is redistributed. Reference adaptation: none was needed and none was made; every reference inside the delivered unit resolves inside it, so no unresolved reference or citation is produced. Printed slips kept literal: no printed word, symbol, hypothesis or number of the counted statement or of its proof was altered. Layout and class: the article's own class is \texttt{amsart} with packages including amssymb, amsxtra, mathtools, hyperref, upref, soul; the wrapper is the lane's pinned \texttt{amsart} editorial wrapper at 10pt on A4, which restates the theorem-like environments the retained text needs and adds the article's own macro definitions that the retained text needs (none), with extra wrapper packages (bm, bbm, dsfont, xspace, cancel, mathtools). The delivered numbering is the unit's own. Assets: no retained span contains a figure environment, a graphics-inclusion command, a PDF-page inclusion, a table or a TikZ picture; no image file is copied into this package, the package declares no asset dependency and no image file is redistributed with it. Licence: version arXiv:2410.13617v2 of this article is distributed under the Creative Commons Attribution 4.0 International licence, as recorded in the arXiv licence metadata for this exact version and shown on the abstract page https://arxiv.org/abs/2410.13617v2. A full rights notice is delivered at licenses/arxiv-2410.13617-CC-BY-4.0.txt. Characters: every retained excerpt is pure ASCII, so the delivered engine, XeLaTeX with its default Latin Modern text font, reports no missing character.
\begin{lemma} \label{lem-stab}
  \leavevmode
  \begin{enumerate}
      \item \label{lem-stab-1} The stabilizer of $\widetilde{Q}_i$ is $W_i$.
      \item \label{lem-stab-2} If $w \widetilde{Q}_i = \widetilde{Q}_j$ where $w \in W$ and $i,j \in [r]$, then $i = j$ and $w \in W_i$.
  \end{enumerate}
\end{lemma}

\begin{lemma} \label{lem-facet}
    \leavevmode
    \begin{enumerate}
      \item \label{lem-facet-1} If the vertex $\lambda$ lies in $w \widetilde{Q}_i$ for some $w \in W$ and some $i \in [r]$, then $\widetilde{Q}_i = w \widetilde{Q}_i$, that is, $w \in W_i$.
      \item \label{lem-facet-2} For any $w \in W$ and $i \in [r]$, if $\widetilde{Q}_i \ne w \widetilde{Q}_i$, then $\widetilde{Q}_i \cap w \widetilde{Q}_i = \emptyset$. 
    \end{enumerate}
\end{lemma}

\begin{lemma} \label{lem-T-coh-gen}
    The algebra $H^*_T(X(\mathit{WP}) ; \mathbb{Q})^W$ is generated by 
    \[\sum_{w \in W^i} X_{i,w}, \quad i \in [r].\]
\end{lemma}

\begin{proof}
    The elements in $H^*_T(X(\mathit{WP}) ; \mathbb{Q})$ are polynomials in the variables $X_{i,w}$.
    Suppose $f \in H^*_T(X(\mathit{WP}) ; \mathbb{Q})^W$ and $f \ne 0$.
    Let $X_{i_1, w_1}^{n_1} \cdots X_{i_k, w_k}^{n_k}$ be a nonzero monomial appearing in $f$, that is,
    \[f = c X_{i_1, w_1}^{n_1} \cdots X_{i_k, w_k}^{n_k} + \text{ other terms, with $c \ne 0$}.\]
    Then by definition, $w_1 \widetilde{Q}_{i_1} \cap \dots \cap w_k \widetilde{Q}_{i_k} \ne \emptyset$.
    There is a vertex of $\mathit{WP}$, say, $w_0\lambda$, lying in this intersection.
    Then, 
    \[\lambda \in w_0^{-1} w_1 \widetilde{Q}_{i_1} \cap \dots \cap w_0^{-1} w_k \widetilde{Q}_{i_k}.\] 
    By Lemma \ref{lem-facet}\eqref{lem-facet-1}, this implies $w_0^{-1} w_j \widetilde{Q}_{i_j} = \widetilde{Q}_{i_j}$ for each $j$. 
    In particular,
    \[w_0^{-1} X_{i_1, w_1}^{n_1} \cdots X_{i_k, w_k}^{n_k} = X_{i_1, e}^{n_1} \cdots X_{i_k, e}^{n_k}.\]
    As $f$ is $W$-invariant, the monomial $X_{i_1, e}^{n_1} \cdots X_{i_k, e}^{n_k}$ also appears in $f$ and we can assume $i_1, \cdots, i_k$ are mutually distinct.

    Let $I := [r] \setminus \{i_1, \dots, i_k\}$, $W_I := \langle s_i \mid i\in I \rangle$ (the parabolic subgroup of $W$ corresponding to $I$), and $W^I$ be the set of minimal representatives for left cosets $W / W_I$.
    Then,
    \[\operatorname{Stab}(X_{i_1, e}^{n_1} \cdots X_{i_k, e}^{n_k}) = \bigcap_{1 \le j \le k} \operatorname{Stab}(X_{i_j,e}) = \bigcap_{1 \le j \le k} W_{i_j} = W_I.\]
    The first equality is due to the fact that $\widetilde{Q}_i$ and $\widetilde{Q}_j$ belong to different $W$-orbit whenever $i \ne j$ (see Lemma \ref{lem-stab}\eqref{lem-stab-2}).
    
    Because $f$ is $W$-invariant, all the $W$-translations of $X_{i_1, e}^{n_1} \cdots X_{i_k, e}^{n_k}$ must appear in $f$ with the same coefficient.
    Therefore, 
    \[f = c \sum_{w \in W^I} X_{i_1, w}^{n_1} \cdots X_{i_k, w}^{n_k} + \text{ other terms, with $c \ne 0$}.\]
    By induction, it suffices to prove the following equality in $H^*_T(X(\mathit{WP}) ; \mathbb{Q})$ :
    \begin{equation} \label{eq-claim-2}
        \sum_{w \in W^I} X_{i_1, w}^{n_1} \cdots X_{i_k, w}^{n_k} = \left( \sum_{w \in W^{i_1}} X_{i_1,w} \right)^{n_1} \cdots \left( \sum_{w \in W^{i_k}} X_{i_k,w} \right)^{n_k}.
    \end{equation}
    By Lemma \ref{lem-stab}\eqref{lem-stab-2}, $w \widetilde{Q}_i$ and $w' \widetilde{Q}_i$ have empty intersection whenever $w \ne w'$.
    Thus, we have $(\sum_{w \in W^i} X_{i,w})^n = \sum_{w \in W^i} X_{i,w}^n$, and Eq.~\eqref{eq-claim-2} is equivalent to 
    \begin{equation} \label{eq-lem-T-coh-gen}
        \sum_{w \in W^I} X_{i_1, w}^{n_1} \cdots X_{i_k, w}^{n_k} =  \sum_{w_1 \in W^{i_1}, \dots, w_k \in W^{i_k}} X_{i_1,w_1}^{n_1}  \cdots  X_{i_k,w_k}^{n_k}.
    \end{equation}
    Note that the summands in the left-hand side of Eq.~\eqref{eq-lem-T-coh-gen} are different to each other, that is, appear only once.
    On the other hand, each summand in the right-hand side of Eq.~\eqref{eq-lem-T-coh-gen}, if nonzero, appears only once as well.
    Therefore, it suffices to show that all the summands in the left-hand side appear in the right-hand side, and vice versa.

    For any $w \in W^I$ and $j \in [k]$, let $w_j$ denote the minimal representative in the coset $wW_{i_j}$.
    Then, the summand $X_{i_1, w}^{n_1} \cdots X_{i_k, w}^{n_k}$ in the left-hand side of Eq.~\eqref{eq-lem-T-coh-gen} equals $X_{i_1, w_1}^{n_1} \cdots X_{i_k, w_k}^{n_k}$, which appears in the right-hand side of Eq.~\eqref{eq-lem-T-coh-gen}.
    Conversely, suppose $X_{i_1, w_1}^{n_1} \cdots X_{i_k, w_k}^{n_k}$ ($w_j \in W^{i_j}$) is a nonzero monomial.
    Then, as above, there is a vertex of $\mathit{WP}$, say, $w' \lambda$, belonging to $w_1 \widetilde{Q}_{i_1} \cap \cdots \cap w_k \widetilde{Q}_{i_k}$.
    Equivalently, we have
    \[\lambda \in (w')^{-1} w_1 \widetilde{Q}_{i_1} \cap \cdots \cap (w')^{-1} w_k \widetilde{Q}_{i_k}.\]
    By Lemma \ref{lem-facet}\eqref{lem-facet-1} again, we have $w_j \in w'W_{i_j}$ for each $j$, and hence $X_{i_j,w_j} = X_{i_j, w'}$.
    Let $w \in W^I$ be the minimal representative in  the coset $w' W_I$.
    Then, 
    \[X_{i_1, w_1}^{n_1} \cdots X_{i_k, w_k}^{n_k} = X_{i_1, w'}^{n_1} \cdots X_{i_k, w'}^{n_k}= X_{i_1, w}^{n_1} \cdots X_{i_k, w}^{n_k},\]
    which is a summand in the left-hand side of Eq.~\eqref{eq-lem-T-coh-gen}.
    The proof is complete.
\end{proof}

\end{document}
